
\begin{abstract}
Over the ten years of the Legacy Survey of Space and Time (LSST), Rubin Observatory will acquire on the order of two million science observations. From twilight to twilight, except during bad weather or maintenance, the Simonyi Telescope will select a new target every 30~seconds, with every observation contributing towards the LSST science goals.  Orchestrating these observations is the work of the Rubin Feature Based Scheduler (FBS).  
The FBS has been generating simulated observations for over ten years, driving the development and co-optimization the LSST survey strategy and FBS algorithms, in a collaborative process between the project and the astronomical community under the guidance of the LSST Survey Cadence Optimization Committee (SCOC). The FBS has been in active use with the Rubin Auxiliary Telescope for several years building the framework for interactions between the FBS and summit system.  The commissioning of LSSTCam in 2025 marked the first on-sky use of the FBS with the Simonyi Telescope.

The commissioning of LSSTCam put the full capability of FBS into action with commissioning configurations similar to those of the full survey configuration. This configuration included multiple modes of observing, which automatically adapt to telemetry about the observatory, night sky and weather conditions during the night, varying from targeted pre-planned Deep Drilling Field observations to wide-area observations designed to acquire pairs of visits every 30 minutes. The FBS also correctly responded to operational constraints such as lunar avoidance zones and configurable altitude or azimuth constraints. During commissioning, we exercised and refined these FBS configurations, adjusting the survey strategy, monitoring FBS performance and survey progress, as well as learning important lessons about the LSST scheduling, the characteristics of resulting observations, and interactions with data management processing. 
\end{abstract}

